% Generated draft: requires mathematical and source review.
\documentclass[11pt]{article}
\usepackage[margin=25mm]{geometry}
\usepackage{fontspec}
\setmainfont{Latin Modern Roman}
\usepackage{amsmath,amssymb,mathtools,unicode-math}
\setmathfont{Latin Modern Math}
\usepackage{enumitem,needspace,xurl,hyperref}
\setlength{\parindent}{0pt}
\setlength{\parskip}{5pt}
\setlength{\emergencystretch}{3em}
\newcommand{\N}{\mathbb{N}}
\newcommand{\Z}{\mathbb{Z}}
\newcommand{\Q}{\mathbb{Q}}
\newcommand{\R}{\mathbb{R}}
\newcommand{\C}{\mathbb{C}}
\newcommand{\F}{\mathbb{F}}
\newcommand{\A}{\mathbb{A}}
\newcommand{\PP}{\mathbb{P}}
\newcommand{\E}{\mathbb{E}}
\newcommand{\eps}{\varepsilon}
\newcommand{\dd}{\,\mathrm{d}}
\begin{document}
\section*{UnsolvedMath ID: 11000021}
\textbf{Source problem code:} catalogue entry
\subsection{Definitions and mathematical statement}
For $g\geq2$, a Hurwitz surface is a compact connected Riemann surface $X$ of genus $g$ attaining the Hurwitz automorphism bound,
\[
 |\operatorname{Aut}_{\mathrm{hol}}(X)|=84(g-1).
\]
In the usual holomorphic-isomorphism convention, put
\[
 h(g)=\#\{[X]:X\text{ is a Hurwitz surface of genus }g\}.
\]
Question 2.20 asks for a formula for $h(g)$ and for the frequency of genera with $h(g)=1$. Genera admitting no Hurwitz surface have $h(g)=0$ and are part of the range unless another sampling convention is specified.

For example, if frequency means natural density among all genera, the quantity to study would be
\[
 \lim_{G\to\infty}
 \frac{\#\{g\in\Z:2\leq g\leq G,\ h(g)=1\}}
 {\#\{g\in\Z:2\leq g\leq G\}}.
\]
This interpretation is conditional: the source does not prescribe natural density, assert that the limit exists, or say whether to condition on $h(g)>0$. It also does not specify the desired kind of formula. Counting up to antiholomorphic equivalence, or counting marked surfaces or group actions instead, would give different counting conventions and is not silently imposed.
\subsection{Short English statement}
Count genus-$g$ Riemann surfaces with the maximal $84(g-1)$ holomorphic automorphisms, and determine how often a genus supports exactly one such surface.
\subsection{Sources}
\begin{itemize}
\item Local registry: \url{data/unsolvedmath/problems.json}, permanent UnsolvedMath ID $11000021$, exact source code \texttt{catalogue entry}.
\item Farb, \emph{Problems on Mapping Class Groups and Related Topics}, Question 2.20, PDF page 30: \url{https://www.math.uchicago.edu/~farb/papers/mcgbook.pdf}. Supplied access date: 2026-07-29; extraction: pdf-font-structure; page-disambiguated numbering among 329 headings.
\item The registry mentions work on Hurwitz/maximal-normal surface counts and quantitative growth without exact titles, authors, or references; it explicitly expresses uncertainty about matching the precise question.
\end{itemize}
\end{document}
